\begin{frame}[label=q5c1]{5c.1 · DataFrame: INNER JOIN geográfico}
  \small
  \pregunta{¿Cómo se distribuye el total entre las regiones definidas por el catálogo geográfico del proyecto?}
  \entrada{Agrupar las denuncias por \texttt{ubigeo}. Unir con \texttt{dimGeografia}, una fila por clave; usar \texttt{region}.}
  \flujo{\texttt{groupBy(ubigeo)}\\Suma de personas}{\texttt{INNER JOIN}\\Clave: ubigeo}{\texttt{groupBy(region)}\\Sumar personas y registros}{3 regiones\\Total 121\,427}
  \begin{columns}[T,onlytextwidth]
    \column{0.55\textwidth}
    \textbf{Código Scala original (fragmento)}
    \codigo{sql/05c_1_inner_region}{15}{18}
    \column{0.42\textwidth}
    \textbf{Salida}\\[5pt]
    {\footnotesize
\begin{tabular}{@{}lll@{}}
\toprule
Región & Personas & Ubigeos\\
\midrule
Costa & 79\,666 & 567\\
Sierra & 34\,301 & 804\\
Selva & 7\,460 & 135\\
\bottomrule
\end{tabular}
    }
  \end{columns}
  \respuesta{Costa reúne 79\,666 personas. El total se conserva porque el catálogo tiene una fila por ubigeo y no hay claves nulas en la salida guardada.}
  \origen{sql}{05c_1_inner_region}
  \note{25 s. Señalar la clave de join y explicar por qué deduplicar la dimensión evita inflar el total.}
\end{frame}
